% Source: https://artofproblemsolving.com/community/c3h2052430_algebra_problem
% Stable identity: aops-topic-2052430-q01
Recall that a quadrilateral $PQRS$ (vertices in this order) is a parallelogram exactly when $\overrightarrow{PQ}=\overrightarrow{SR}$, i.e.\ $Q-P=R-S$, which is equivalent to
\[P+R=Q+S .\tag{1}\]
In words: the sums of opposite vertices are equal, i.e.\ the two diagonals $PR$ and $QS$ have the same midpoint.

\begin{asy}
usepackage("amsmath");
size(6cm);
defaultpen(fontsize(10pt));
pair w=(0,0), y=(4,0.6), x=(5.2,2.8), z=(1.2,2.2);
draw(w--y--x--z--cycle,linewidth(0.9));
draw(w--x,dashed);draw(y--z,dashed);
pair Mid=(w+x)/2;
dot(w);dot(x);dot(y);dot(z);dot(Mid,red);
label("$w$",w,SW);label("$y$",y,SE);label("$x$",x,NE);label("$z$",z,NW);
label("$O$",Mid,N,red);
label("$O=\dfrac{w+x}{2}=\dfrac{y+z}{2}$",(2.6,-0.9),red);
label("\mbox{}",(min(currentpicture,user=true)+max(currentpicture,user=true))/2,black); // keeps the text after the figure black when LaTeX typesets the labels inline
\end{asy}

\textbf{Only if.} If the four points are the vertices of a parallelogram $PQRS$, then (1) holds, so the sum of the two opposite vertices $P,R$ equals the sum of the other two.

\textbf{If.} Suppose, after renaming, $w+x=y+z$. Then $y-w=x-z$, so the vectors from $w$ to $y$ and from $z$ to $x$ are equal and nonzero (the points are distinct). The segments $wy$ and $zx$ are therefore parallel and of equal length. They do not lie on one line, since otherwise $w,y,z$ would be collinear. Hence $w,y,x,z$, taken in this order, is a quadrilateral with a pair of opposite sides $wy$ and $zx$ equal and parallel, so it is a parallelogram, with diagonals $wx$ and $yz$ meeting at their common midpoint $\dfrac{w+x}{2}=\dfrac{y+z}{2}$.

\textbf{The hypothesis is needed.} The points $0,1,2,3$ on the real axis satisfy $0+3=1+2$, but they are collinear and are not the vertices of any parallelogram. $\blacksquare$
